\documentclass{article}
\usepackage{amsmath, amssymb, amsfonts}
\usepackage[utf8]{inputenc}
\usepackage[spanish]{babel}

\title{Complejos}
\author{Iván Palomo, Alexandre Torrandell}
\date{23 de febrero de 2026}

\begin{document}
	\maketitle
	\section{Exercici 2 complexos}
		
Demostrad que si $z_1$ i $z_2$ son raices del polinomio con coeficientes reales $x^2+ax+b$, aleshores $$(z_1)^n+(z_2)^\in R$$.
Por tanto al ser una ecuacion cuadratica podemos tener dos casos. Que $z_1$ sea real hace que $z_2$ tambien lo sea y por tanto la suma de ellos elevados a n es trivial. En el otro caso si $z_1$ es complejo al ser cuadratica $z_2$ es compleja i ademas el complementario de este. 
	
Para facilitar la resolución consideramos $z_1$ y $z_2$ como forma polar. Donde $z_2$ tendra el angulo teta negativo (por ser el conjugado). Si aplicamos la formula que pasa de polar a funciones trigonometricas aplicado a nombres complejos elevados veremos que se cancela el termino $i\sin(n\theta)$. Ya que en el $\cos$ el aungulo negativo y positivo es el mismo. Por otro lado en el $\sin$ el signo sale afuera permitiendo que se resten la parte imaginaria de la expresion.



\end{document}